% =============================================================================
%  ELEMENTOS PRE-TEXTUAIS (adaptados para classe book)
% =============================================================================

% --- AGRADECIMENTOS ---
\chapter*{Agradecimentos}
\label{ch:agradecimentos}
\markboth{AGRADECIMENTOS}{}

Este livro é fruto de anos de pesquisa, experimentação e desenvolvimento 
contínuo do \opencodeeco{}. Agradeço profundamente à comunidade de 
engenharia de software e inteligência artificial que, através de décadas 
de pesquisa aberta e colaborativa, tornou possível a construção de 
sistemas cognitivos artificiais cada vez mais sofisticados.

Agradecimentos especiais aos pesquisadores e engenheiros cujo trabalho 
fundamenta cada capítulo desta obra: John McCarthy, Marvin Minsky, 
Alan Turing, Claude Shannon, John von Neumann, Donald Knuth, Kent Beck, 
Martin Fowler, Stuart Russell, Peter Norvig, Yoshua Bengio, Geoffrey 
Hinton, Yann LeCun, Andrew Ng, e tantos outros que pavimentaram o 
caminho da inteligência artificial e da engenharia de software.

À Universidade Federal do Ceará (UFC), campus Crateús, pelo ambiente 
acadêmico que estimula a pesquisa interdisciplinar. Ao Programa de 
Pós-Graduação em Tecnologia e Engenharia (PPGTE), pela estrutura 
que permitiu a validação científica dos conceitos aqui apresentados.

Aos estudantes e pesquisadores que testaram, questionaram e aprimoraram 
cada componente do \opencodeeco{} ao longo de 23 ciclos evolutivos 
(R1 a R23), meu mais sincero reconhecimento.

Finalmente, à minha família, pelo apoio incondicional durante os 
longos períodos de imersão necessários à conclusão desta obra.

% --- DEDICATORIA ---
\chapter*{Dedicatória}
\label{ch:dedicatoria}
\markboth{DEDICATÓRIA}{}

\vspace*{\fill}
\begin{flushright}
\textit{A todos os estudantes que ousam questionar,\\
aos engenheiros que constroem o futuro\\
e aos pesquisadores que jamais deixam de aprender.}
\end{flushright}
\vspace*{\fill}

% --- EPIGRAFE ---
\chapter*{Epígrafe}
\label{ch:epigrafe}
\markboth{EPÍGRAFE}{}

\vspace*{\fill}
\begin{flushright}
\textit{``The most profound technologies are those that disappear.}\\
\textit{They weave themselves into the fabric of everyday life}\\
\textit{until they are indistinguishable from it.''}\\[0.3cm]
--- Mark Weiser, \textit{The Computer for the 21st Century}, 1991\par
\vspace{1cm}
\textit{``All models are wrong, but some are useful.''}\\[0.3cm]
--- George E. P. Box, \textit{Science and Statistics}, 1976\par
\vspace{1cm}
\textit{``The best way to predict the future is to invent it.''}\\[0.3cm]
--- Alan Kay, \textit{PARC}, 1971\par
\end{flushright}
\vspace*{\fill}

% --- RESUMO ---
\chapter*{Resumo}
\label{ch:resumo}
\markboth{RESUMO}{}

Este livro apresenta uma jornada autodidata e sistemática através dos fundamentos 
matemáticos, estatísticos e computacionais que sustentam a engenharia de 
ecossistemas cognitivos artificiais, tendo como objeto central de estudo e 
experimentação o \opencodeeco{} (versão 5.4.0, ciclo evolutivo R23). 
A obra estrutura-se em oito capítulos que conduzem o leitor do nível zero 
(fundamentos básicos de matemática e lógica) ao nível PhD (validação 
científica, produção acadêmica \qualisA{} e defesa perante banca examinadora).

O Capítulo 1 estabelece a base matemática e estatística --- álgebra linear, 
cálculo, probabilidade, inferência estatística e teoria da informação --- 
essencial para a compreensão dos algoritmos de inteligência artificial. 
O Capítulo 2 introduz os conceitos fundamentais de IA, aprendizado de máquina, 
redes neurais profundas (transformers) e arquiteturas de agentes autônomos. 
O Capítulo 3 mergulha na arquitetura do \opencodeeco, detalhando o padrão 
três camadas (MCP $\rightarrow$ Skill $\rightarrow$ Agent), a metodologia 
SDD+TDD, o barramento de eventos e o sistema de injeção de dependência.

O Capítulo 4 explora o \textit{Pipeline de Scanners Epistemológicos} --- 
Noológico, Teleológico, Evolutivo, Refinamento e MCSP --- e o sistema de 
metacognição (SPEC-036) com Self-Model N0-N3 e monitor metacognitivo. 
O Capítulo 5 apresenta o Trust Engine (SPEC-038) com Behavioral Gate, 
Natural Forgetting (modelo Atkinson-Shiffrin) e Governança Cooperativa 
(princípios de Ostrom). O Capítulo 6 detalha a Token Economy (SPEC-022/023/024) 
com \textit{staking}, \textit{slashing}, mercado de taxas e trilha de auditoria.

O Capítulo 7 descreve o sistema de experimentação e validação científica: 
benchmark CORA-Eval (150 tarefas, 10 dimensões), validação Aletheia 
(Lean 4), pipeline de produção acadêmica MASWOS com 49 agentes especializados 
e o sistema de correção iterativa Qualis A1. O Capítulo 8, finalmente, 
apresenta o roteiro completo para produção e defesa de dissertação: 
metodologia PPGTE/UFC, simulação de banca com agentes, protocolo de 
anonimato e métricas de avaliação DAP.

Cada capítulo contém definições formais, ilustrações (TikZ), exemplos 
práticos executáveis no \opencodeeco{}, exercícios progressivos (nível 0 
ao PhD), notas de rodapé com citações verificáveis e links ativos para 
fontes primárias.

\medskip
\noindent\textbf{Palavras-chave}: Engenharia de Software. Inteligência Artificial. 
Sistemas Multiagentes. Metacognição. Trust Engine. OpenCode Ecosystem. 
SDD+TDD. Qualis A1.

% --- ABSTRACT ---
\chapter*{Abstract}
\label{ch:abstract}
\markboth{ABSTRACT}{}

\begin{otherlanguage*}{english}

This book presents a self-contained and systematic journey through the 
mathematical, statistical, and computational foundations underlying the 
engineering of artificial cognitive ecosystems, with the \opencodeeco{} 
(v5.4.0, evolutionary cycle R23) as the central object of study and 
experimentation. The work is structured in eight chapters that guide the 
reader from level zero (basic mathematics and logic) to PhD level 
(scientific validation, \qualisA{} academic production, and defense 
before an examining committee).

Chapter 1 establishes the mathematical and statistical foundations --- 
linear algebra, calculus, probability, statistical inference, and 
information theory --- essential for understanding artificial intelligence 
algorithms. Chapter 2 introduces fundamental AI concepts, machine learning, 
deep neural networks (transformers), and autonomous agent architectures. 
Chapter 3 delves into the \opencodeeco{} architecture, detailing the 
three-layer pattern (MCP $\rightarrow$ Skill $\rightarrow$ Agent), the 
SDD+TDD methodology, the event bus, and the dependency injection system.

Chapter 4 explores the Epistemological Scanner Pipeline --- Noological, 
Teleological, Evolutionary, Refinement, and MCSP --- and the metacognition 
system (SPEC-036) with Self-Model N0-N3 and metacognitive monitoring. 
Chapter 5 presents the Trust Engine (SPEC-038) with Behavioral Gate, 
Natural Forgetting (Atkinson-Shiffrin model), and Cooperative Governance 
(Ostrom's principles). Chapter 6 details the Token Economy 
(SPEC-022/023/024) with staking, slashing, fee market, and audit trail.

Chapter 7 describes the experimentation and scientific validation system: 
CORA-Eval benchmark (150 tasks, 10 dimensions), Aletheia validation 
(Lean 4), the MASWOS academic production pipeline with 49 specialized 
agents, and the iterative Qualis A1 correction system. Chapter 8, 
finally, presents a complete roadmap for dissertation production and 
defense: PPGTE/UFC methodology, board examination simulation with agents, 
anonymity protocol, and DAP evaluation metrics.

Each chapter contains formal definitions, illustrations (TikZ), practical 
examples executable within the \opencodeeco{}, progressive exercises 
(levels 0 to PhD), footnotes with verifiable citations, and active links 
to primary sources.

\medskip
\noindent\textbf{Keywords}: Software Engineering. Artificial Intelligence. 
Multiagent Systems. Metacognition. Trust Engine. OpenCode Ecosystem. 
SDD+TDD. Qualis A1.

\end{otherlanguage*}
